
\subsection*{Dataset construction}

We constructed a corpus of Health Informatics publications from PubMed by selecting records indexed with the MeSH Major Topic ``Health Informatics'' (MeSH ID D008490). We restricted the retrieval to journal articles and excluded reviews, systematic reviews, meta-analyses, editorials, letters, comments, news items, biographies, errata, and retracted publications. The initial retrieval yielded 86,549 unique records within the study windows. We excluded 6,449 records without abstracts and 10,891 records that were not journal articles, leaving a final analytic sample of 69,209 publications.

To examine temporal shifts associated with the emergence of LLMs, we partitioned the corpus into two periods anchored to the public release of ChatGPT-3.5\cite{openai2022chatgpt} on November 30, 2022. 
The pre-LLM period contained publications from January 1, 2015, through November 30, 2022 ($n=53,575$), while the post-LLM period contained publications from December 1, 2022, to February 1, 2026 ($n=15,634$).

\subsection*{Identification of LLM-assisted writing}

We developed a supervised classifier to identify LLM-assisted scientific writing.
A SciBERT model\cite{beltagy-etal-2019-scibert}  was fine-tuned on a curated training set of original pre-LLM abstracts and LLM-generated rewrites. The training set consisted of 2,500 original abstracts and 2,500 corresponding GPT-4.1\cite{openai2025gpt41} generated versions produced using four rewriting strategies: full rewrite (Supplementary Prompt~1), polish (Supplementary Prompt~2), light edit (Supplementary Prompt~3), and human-like rewrite (Supplementary Prompt~4). Original abstracts and their rewritten counterparts were assigned to the same data split to prevent information leakage.

The classifier was calibrated to prioritize precision over recall. For threshold selection, we constructed  a balanced labeled dataset of 14,368 abstracts and five pre-LLM negative-control samples of 2,500 abstracts each, excluding the abstracts used to construct the training pairs. At the selected threshold of approximately 0.999876, precision was 0.981, recall was 0.441, and F1 was 0.608 on the labeled threshold-tuning dataset. Across the pre-LLM samples, the mean false-positive rate was 0.00912 and the maximum was 0.010. 
This conservative threshold was selected because the aim of the study was not to identify all LLM-assisted publications, but rather to define a high-confidence subset of papers with likely LLM-associated writing. 

Publications in the post-LLM period with classifier scores above this threshold were assigned to the LLM-assisted group ($n = 2,673$). The remaining post-LLM publications were assigned to the LLM-independent group ($n = 12,961$). This design allowed us to distinguish broader temporal shifts from additional effects associated with LLM-assisted writing within the post-LLM era.

\subsection*{Feature extraction}
We characterized publications in multiple complementary features because no single one can simultaneously capture semantic content, focal subject matter, and cited knowledge inputs (Supplementary Table~S1). These features support different downstream analyses, including semantic concentration, topic composition, and interdisciplinarity.


\subsubsection*{Semantic representation}
To capture semantic content across publications, we encoded each title and abstract with SPECTER2\cite{singh2022scirepeval,allenai2023specter2base}. This model provides dense representations of scientific documents in a shared embedding space. Following the official implementation guidelines, we used the base model together with its proximity adapter to encode each publication into a 768-dimensional vector.

\subsubsection*{Focal topics}
We used MeSH descriptors assigned in PubMed to characterize focal topics\cite{NLMMeSH2026}. MeSH descriptors and OpenAlex topic labels were mapped to a shared ontology with 31 subtopics under Biological Sciences and Health Sciences. Biological Sciences subtopics followed the Nature Portfolio subject taxonomy, while Health Sciences subtopics were consolidated to capture clinically oriented themes relevant to health informatics. Azure OpenAI GPT-4.1, version 2025-04-14, assigned an ontology subtopic, confidence score, and rationale to each label using the system and user prompts provided in Supplementary Prompts~5 and~6. The mapping file contained 13,540 MeSH descriptors and 2,565 OpenAlex labels. Of these 16,105 labels, 14,765 had confidence scores of at least 0.70 and 1,340 had lower scores. The review flag covered 1,347 labels, including the 1,340 lower-confidence labels and seven higher-confidence labels. These flags identify assignments selected for inspection and do not themselves document completion of manual review.

\subsubsection*{Reference topics}
To characterize the topical breadth of cited knowledge, we retrieved each referenced paper's OpenAlex topic assignment\cite{priem2022openalex}, mapped those labels into the same 31-subtopic ontology, and used the resulting reference-topic profiles to quantify the topical composition of each bibliography and drive the interdisciplinary measures.


\subsubsection*{Language barrier}

English remains the dominant language of indexed scientific publishing, including most peer-reviewed journals and international conferences, potentially disadvantaging researchers whose primary language is not English. Here, we investigate whether LLMs help mitigate such language-related barriers.
We define a paper \textit{Anglophone} if the first author's affiliation country appeared in either of two Wikipedia categories, ``English is a de jure official language'' or ``English is a predominant language without de jure designation as official'' \cite{wikipediaContributors2026list}. 

The countries are: 
\begin{enumerate*}[label=(\arabic*)]
\item {Antigua and Barbuda},
\item {Australia},
\item Bahamas,
\item Barbados,
\item {Belize},
\item {Botswana},
\item {Burundi},
\item {Cameroon},
\item {Canada},
\item Dominica,
\item {Eswatini},
\item {Fiji},
\item {Gambia},
\item {Ghana},
\item {Grenada},
\item {Guyana},
\item {India},
\item {Ireland},
\item {Jamaica},
\item {Kenya},
\item Kiribati,
\item {Lesotho},
\item {Liberia},
\item {Malawi},
\item {Malta},
\item Marshall Islands,
\item Micronesia,
\item {Namibia},
\item Nauru,
\item {New Zealand},
\item {Nigeria},
\item {Pakistan},
\item Palau,
\item {Papua New Guinea},
\item {Philippines},
\item {Rwanda},
\item {Saint Kitts and Nevis},
\item Saint Lucia,
\item Saint Vincent and the Grenadines,
\item Samoa,
\item {Seychelles},
\item {Sierra Leone},
\item {Singapore},
\item Solomon Islands,
\item {South Africa},
\item {South Sudan},
\item {Sudan},
\item {Tanzania},
\item Tonga,
\item {Trinidad and Tobago},
\item Tuvalu,
\item {Uganda},
\item {United Kingdom},
\item {United States},
\item Vanuatu,
\item {Zambia},
\item {Zimbabwe}.
\end{enumerate*}


\subsubsection*{Global access gap.}

Scholars in low- and middle-income economies (LMICs) have historically faced barriers to scientific publishing, including limited access to literature and resource constraints. We examine whether LLM tools are associated with improvements in research output from these settings.
We define a paper as \textit{LMIC} if the last author is affiliated with institutions in low- and middle-income economies based on World Bank Country Classifications.  

Low-income economies include:
\begin{enumerate*}[label=(\arabic*)]
\item Afghanistan,
\item Burkina Faso,
\item Burundi,
\item Central African Republic,
\item Chad,
\item Congo, Dem. Rep.,
\item Eritrea,
\item Gambia, The,
\item Guinea-Bissau,
\item Korea, Dem. People's Rep.,
\item Liberia,
\item Madagascar,
\item Malawi,
\item Mali,
\item Mozambique,
\item Niger,
\item Rwanda,
\item Sierra Leone,
\item Somalia, Fed. Rep.,
\item South Sudan,
\item Sudan,
\item Syrian Arab Republic,
\item Togo,
\item Uganda,
\item Yemen, Rep.
\end{enumerate*}

Lower middle income economies include:
\begin{enumerate*}[label=(\arabic*)]
\item Angola,
\item Bangladesh,
\item Benin,
\item Bhutan,
\item Bolivia,
\item Cambodia,
\item Cameroon,
\item Comoros,
\item Congo, Rep.,
\item Côte d’Ivoire,
\item Djibouti,
\item Egypt, Arab Rep.,
\item Eswatini,
\item Ghana,
\item Guinea,
\item Haiti,
\item Honduras,
\item India,
\item Jordan,
\item Kenya,
\item Kiribati,
\item Kyrgyz Republic,
\item Lao PDR,
\item Lebanon,
\item Lesotho,
\item Mauritania,
\item Micronesia, Fed. Sts.,
\item Morocco,
\item Myanmar,
\item Namibia,
\item Nepal,
\item Nicaragua,
\item Nigeria,
\item Pakistan,
\item Papua New Guinea,
\item Philippines,
\item São Tomé and Príncipe,
\item Senegal,
\item Solomon Islands,
\item Sri Lanka,
\item Tajikistan,
\item Tanzania,
\item Timor-Leste,
\item Tunisia,
\item Uzbekistan,
\item Vanuatu,
\item Vietnam,
\item West Bank and Gaza,
\item Zambia,
\item Zimbabwe.
\end{enumerate*}

\subsubsection*{International collaboration.}

To characterize collaboration structure, we defined \textit{international collaboration} as papers with author affiliations spanning more than one continent.


\subsection*{Outcome measures}


We evaluated publications using bibliometric, semantic, topical and authorship-related measures.
In addition to standard variables such as raw cited-by count, journal impact factor, and number of references, we introduced several derived measures for this study. These measures were designed to make papers comparable across publication windows, summarize the topical structure of their cited and focal knowledge, quantify semantic concentration in embedding space, and characterize authorship representation across country contexts.

\paragraph{Distance to topic centroid.}
We computed the Euclidean distance between a paper's embedding and the centroid of its focal topic in the whole cohort (pre-LLM, LLM-independent, LLM-assisted).

\paragraph{Reference-list length.} Reference-list length was measured using the number of references recorded in the saved OpenAlex-derived per-paper data. It measures publications cited by a paper, while cumulative citations measure subsequent publications citing that paper.

\paragraph{Rao-Stirling diversity of reference topics.}
To measure the breadth and diversity of cited knowledge, we computed \textit{Rao-Stirling diversity} of each paper's reference topic profile, defined as $D_p = \sum_i \sum_j p_i p_j d_{ij}$, where $p_i$ is the proportion of paper $p$'s references assigned to topic $i$, and $d_{ij}$ is the disparity between topics $i$ and $j$. We set $d_{ij}=0$ for identical topics, $d_{ij}=0.5$ for different subtopics within the same major ontology domain, and $d_{ij}=1$ for subtopics from different major domains.

\paragraph{Focal-reference Rao-Stirling distance.}
Similarly, to measure how far a paper's cited knowledge extends from its own focal topic profile, we computed \textit{focal-reference Rao-Stirling distance}, defined as $D_p^{FR} = \sum_i \sum_j f_i r_j d_{ij}$, where $f_i$ is the proportion of paper $p$'s focal topics assigned to topic $i$, $r_j$ is the proportion of paper $p$'s references assigned to topic $j$, and $d_{ij}$ is the disparity between focal and reference topics. We used the same ontology distance definition: $d_{ij}=0$ for identical topics, $d_{ij}=0.5$ for different subtopics within the same major ontology domain, and $d_{ij}=1$ for subtopics from different major domains.

\paragraph{Cumulative citation acquisition.}
To account for differences in citation exposure time across publication dates, we retrieved exact citing works from OpenAlex and counted non-negative-lag citation events by the citing work's publication date. For each paper $p$ and post-publication month $m$, we computed cumulative citations as $C_p(m)$, the number of citations received from publication through month $m$. A paper contributed to month $m$ only if that month was observable by the analysis date. The primary visibility analysis used cumulative citations through 24 months after publication; citations occurring before the target paper's publication date were excluded.
Cumulative citation comparisons at 12, 24, and 36 months are reported in Supplementary Table~S2. Each window includes only papers with sufficient observable follow-up. These comparisons use different eligible subsets and should not be read as changes within a fixed sample. No LLM-assisted papers had sufficient follow-up for a 48-month comparison.

\paragraph{Journal impact factor.} We used the OpenAlex source-level field \nolinkurl{summary_stats.2yr_mean_citedness} as the journal impact factor metric in this study. Available values were summarized within the validated cohort.

\paragraph{Journal access model.} We used OpenAlex journal-level access metadata to classify fully open-access journals, including fee-free diamond journals, as Gold Open Access and other journals as subscription/hybrid. Open-access articles in hybrid journals remained in the subscription/hybrid category. Papers without usable metadata and records assigned to non-journal sources were excluded. Supplementary Table~S3 reports the proportions among papers with identifiable status. Sensitivity analyses used Directory of Open Access Journals membership and known transitions to fully open access, excluding papers published in the transition year.

\subsection*{Statistical analysis}

Because LLM-assisted papers were necessarily published in the post-LLM period, we separated the analysis into two contrasts: LLM-assisted versus post-era LLM-independent papers, and LLM-assisted papers versus pre-LLM papers. This design allowed us to distinguish patterns associated with LLM-assisted writing from broader temporal differences between the pre- and post-LLM periods. For outcomes shown with statistical annotations, we used two-sided Mann--Whitney U tests\cite{mann1947test}. Journal-access distributions were compared using chi-square tests, with Yates correction for pairwise two-by-two comparisons. We adjusted the three pairwise comparisons using the Benjamini--Hochberg procedure. The cumulative citation table reports unadjusted two-sided Mann--Whitney U test values alongside the corresponding adjusted values.

